% Part: incompleteness
% Chapter: incompleteness-provability
% Section: rosser-thm

\documentclass[../../../include/open-logic-section]{subfiles}

\begin{document}

\olfileid{inc}{inp}{ros}

\olsection{مبرهنة روسر}

هل يمكن تقوية برهان غودل بإسقاط قيد الاتساق-$\omega$ والاكتفاء
بالاتساق؟ نعم؛ وهذه هي حيلة روسر: يُستبدل بمحمول
!!{derivability} المعتاد محمول معدّل، هو $\ORProv_\OLId{latin-looped}{T}(\OLId{latin-ordinary}{y})$،
مكان $\OProv[\Th{T}](\OLId{latin-ordinary}{y})$.

\begin{thm}
\ollabel{thm:rosser}
إذا كانت $\Th{T}$ نظرية متسقة و!!{axiomatizable} تمتد $\Th{Q}$،
فإن $\Th{T}$ غير تامة.
\end{thm}

\begin{proof}
تعريف $\OProv[\Th{T}](\OLId{latin-ordinary}{y})$ المتقدم هو
$\lexists[\OLId{latin-ordinary}{x}][\OPrf[\Th{T}](\OLId{latin-ordinary}{x},\OLId{latin-ordinary}{y})]$، والصيغة
$\OPrf[\Th{T}](\OLId{latin-ordinary}{x},\OLId{latin-ordinary}{y})$ تمثل العلاقة القابلة للقرار التي معناها أن
$\OLId{latin-ordinary}{x}$ عدد غودل لـ!!a{derivation} للـ!!{sentence} ذات عدد غودل $\OLId{latin-ordinary}{y}$.
وكذلك تقبل القرار العلاقة بين $\OLId{latin-ordinary}{x}$ و$\OLId{latin-ordinary}{y}$ التي يكون بها $\OLId{latin-ordinary}{x}$ عدد غودل
لـ\emph{دحض} الجملة ذات عدد غودل $\OLId{latin-ordinary}{y}$. فنأخذ $\fn{not}(\OLId{latin-ordinary}{x})$
الدالة العودية البدائية التالية: إذا رمز $\OLId{latin-ordinary}{x}$ الصيغة $!A$،
أعطت $\fn{not}(\OLId{latin-ordinary}{x})$ شفرة $\lnot !A$. فصدق
$\Refut[\Th{T}](\OLId{latin-ordinary}{x},\OLId{latin-ordinary}{y})$ يكافئ صدق
$\Prf[\Th{T}](\OLId{latin-ordinary}{x},\fn{not}(\OLId{latin-ordinary}{y}))$؛ ولتكن
$\ORefut[\Th{T}](\OLId{latin-ordinary}{x},\OLId{latin-ordinary}{y})$ الصيغة الممثلة لهذه العلاقة.
فإذا كان $\Th{T}\Proves\lnot !A$، وكانت $\delta$
!!{derivation} لذلك، لزم
$\Th{Q}\Proves\ORefut[\Th{T}](\gn{\delta},\gn{!A})$.
ونضع $\ORProv[\Th{T}](\OLId{latin-ordinary}{y})$ للتعبير
\[
\lexists[\OLId{latin-ordinary}{x}][(\OPrf[\Th{T}](\OLId{latin-ordinary}{x},\OLId{latin-ordinary}{y}) \land \lforall[\OLId{latin-ordinary}{z}][(\OLId{latin-ordinary}{z} < \OLId{latin-ordinary}{x} \lif \lnot
  \ORefut[\Th{T}](\OLId{latin-ordinary}{z},\OLId{latin-ordinary}{y}))])].
\]
فالمعنى التقريبي لـ$\ORProv[\Th{T}](\OLId{latin-ordinary}{y})$ هو: «يوجد برهان على $\OLId{latin-ordinary}{y}$
في $\Th{T}$، وليس ثمة دحض لـ$\OLId{latin-ordinary}{y}$ أصغر منه في عدد غودل».
ومتى كانت $\Th{T}$ متسقة، صدقت $\ORProv[\Th{T}](\OLId{latin-ordinary}{y})$ على
الأعداد عينها التي تصدق عليها $\OProv[\Th{T}](\OLId{latin-ordinary}{y})$.
غير أن للمحمولين خصائص مختلفة من جهة \emph{قابلية الإثبات}
في $\Th{T}$؛ وقد تبيّن لنا أن الصدق غير قابلية الإثبات.
أما إذا كانت $\Th{T}$ \emph{غير} متسقة، فهما \emph{لا} يتفقان
حتى في الأعداد التي يصدقان عليها!{} ويشيع أن تُقرأ
$\ORProv[\Th{T}](\OLId{latin-ordinary}{y})$ بقولنا «$\OLId{latin-ordinary}{y}$ قابلة للإثبات وفق روسر».
وقد توهم هذه القراءة أن المقصود نوع خاص من الإثبات المعتاد،
مع أن في النظريات غير المتسقة !!{sentence}s تثبت ولا تثبت وفق
روسر. ودفعًا لهذا اللبس يمكن اختيار اسم مصنوع لا يوهم ذلك،
كقولنا «$\OLId{latin-ordinary}{y}$ قابلة للشَّمْوَلة».

وبلمّة النقطة الثابتة توجد صيغة $!R_\Th{T}$ تحقق
\begin{equation}
  \Th{Q} \Proves !R_\Th{T} \liff \lnot \ORProv[\Th{T}](\gn{!R_\Th{T}}).
  \ollabel{RT}
\end{equation}
ونخالف هنا برهان \olref[1in]{thm:first-incompleteness}:
يكفي اتساق $\Th{T}$ لكي يمتنع على $\Th{T}$ !!{derive} الصيغة
$!R_\Th{T}$، ويمتنع على $\Th{T}$ أيضًا !!{derive} نفيها
$\lnot !R_\Th{T}$. فلا حاجة إلى الاتساق-$\omega$.

نثبت أولًا $\Th{T}\Proves/ !R_{\Th{T}}$. لو ثبتت، لوجد
!!a{derivation} لـ$!R_{\Th{T}}$ من $\Th{T}$؛ وليكن $\OLId{latin-ordinary}{n}$
عدد غودل له. عندئذ
$\Th{Q}\Proves\OPrf[\Th{T}](\num{\OLId{latin-ordinary}{n}},\gn{!R_{\Th{T}}})$،
إذ $\OPrf[\Th{T}]$ تمثل $\Prf[\Th{T}]$ في $\Th{Q}$.
ولكل $\OLId{latin-ordinary}{k}<\OLId{latin-ordinary}{n}$ لا يكون $\OLId{latin-ordinary}{k}$ عدد غودل لـ!!a{derivation}
لـ$\lnot !R_{\Th{T}}$، لاتساق $\Th{T}$.
فلكل $\OLId{latin-ordinary}{k}<\OLId{latin-ordinary}{n}$:
$\Th{Q}\Proves\lnot\ORefut[\Th{T}](\num{\OLId{latin-ordinary}{k}},\gn{!R_{\Th{T}}})$.
وتعطينا \olref[req][min]{lem:less-nsucc}
$\Th{Q}\Proves\lforall[\OLId{latin-ordinary}{z}][(\OLId{latin-ordinary}{z}<\num{\OLId{latin-ordinary}{n}}\lif
\lnot\ORefut[\Th{T}](\OLId{latin-ordinary}{z},\gn{!R_{\Th{T}}}))]$. ومن ثم
\[
\Th{Q} \Proves \lexists[\OLId{latin-ordinary}{x}][(\OPrf[\Th{T}](\OLId{latin-ordinary}{x},\gn{!R_{\Th{T}}}) \land \lforall[\OLId{latin-ordinary}{z}][(\OLId{latin-ordinary}{z}
    < \OLId{latin-ordinary}{x} \lif \lnot \ORefut[\Th{T}](\OLId{latin-ordinary}{z},\gn{!R_{\Th{T}}}))])],
\]
وهذا هو $\ORProv[\Th{T}](\gn{!R_{\Th{T}}})$.
فبمقتضى \olref{RT} يلزم $\Th{Q}\Proves\lnot !R_{\Th{T}}$.
ولأن $\Th{T}$ تمتد $\Th{Q}$، يلزم
$\Th{T}\Proves\lnot !R_{\Th{T}}$ أيضًا.
وقد فرضنا $\Th{T}\Proves !R_{\Th{T}}$، فتكون $\Th{T}$
غير متسقة، خلاف المطلوب في فرض المبرهنة.

ونثبت ثانيًا $\Th{T}\Proves/\lnot !R_{\Th{T}}$.
لو ثبت النفي، فليكن $\OLId{latin-ordinary}{n}$ عدد غودل لـ!!a{derivation}
لـ$\lnot !R_{\Th{T}}$. فحينئذ
$\Refut[\Th{T}](\OLId{latin-ordinary}{n},\Gn{!R_{\Th{T}}})$ صادقة؛
وبما أن $\ORefut[\Th{T}]$ تمثل $\Refut[\Th{T}]$ في $\Th{Q}$،
يلزم $\Th{Q}\Proves\ORefut[\Th{T}](\num{\OLId{latin-ordinary}{n}},\gn{!R_{\Th{T}}})$.
ونبلغ التناقض مرة أخرى بإثبات أن $\Th{T}$ يمكنها !!{derive}
الصيغة $!R_{\Th{T}}$ أيضًا. فقد علمنا أن
\begin{align*}
\Th{Q} & \Proves !R_{\Th{T}} \liff \lnot \ORProv[\Th{T}](\gn{!R_{\Th{T}}}),
\intertext{ولما كانت $\Th{T}$ تمتدّ~$\Th{Q}$، فيكفي أن نبيّن أن}
\Th{Q} & \Proves \lnot \ORProv[\Th{T}](\gn{!R_{\Th{T}}}).
\end{align*}
والـ!!{sentence} $\lnot\ORProv[\Th{T}](\gn{!R_{\Th{T}}})$،
أي
\begin{align*}
  \lnot & \lexists[\OLId{latin-ordinary}{x}][(\OPrf[\Th{T}](\OLId{latin-ordinary}{x},\gn{!R_{\Th{T}}}) \land
    \lforall[\OLId{latin-ordinary}{z}][(\OLId{latin-ordinary}{z} < \OLId{latin-ordinary}{x} \lif \lnot \ORefut[\Th{T}](\OLId{latin-ordinary}{z},\gn{!R_{\Th{T}}}))])],
  \intertext{مكافئة منطقيًا لـ}
  & \lforall[\OLId{latin-ordinary}{x}][(\OPrf[\Th{T}](\OLId{latin-ordinary}{x},\gn{!R_{\Th{T}}}) \lif
    \lexists[\OLId{latin-ordinary}{z}][(\OLId{latin-ordinary}{z} < \OLId{latin-ordinary}{x} \land \ORefut[\Th{T}](\OLId{latin-ordinary}{z},\gn{!R_{\Th{T}}}))])].
\end{align*}
نستدل الآن استدلالًا غير صوري بالمنطق وبما !!{derive}s النظام
$\Th{Q}$. نأخذ $\OLId{latin-ordinary}{x}$ اعتباطيًا ونفرض
$\OPrf[\Th{T}](\OLId{latin-ordinary}{x},\gn{!R_{\Th{T}}})$.
وبما أنه قد ثبت $\Th{T}\Proves/ !R_{\Th{T}}$،
فلكل $\OLId{latin-ordinary}{k}$:
$\Th{Q}\Proves\lnot\OPrf[\Th{T}](\num{\OLId{latin-ordinary}{k}},\gn{!R_{\Th{T}}})$.
إذن، لكل $\OLId{latin-ordinary}{k}$، يلزم $\eq/[\OLId{latin-ordinary}{x}][\num{\OLId{latin-ordinary}{k}}]$؛ ومنه بخاصة
(أ) $\eq/[\OLId{latin-ordinary}{x}][\num{\OLId{latin-ordinary}{n}}]$. ويلزم أيضًا
$\lnot(\eq[\OLId{latin-ordinary}{x}][\num{\OLMathNumeral{0}}]\lor\eq[\OLId{latin-ordinary}{x}][\num{\OLMathNumeral{1}}]\lor\dots\lor
\eq[\OLId{latin-ordinary}{x}][\num{\OLId{latin-ordinary}{n}-\OLMathNumeral{1}}])$.
فمن \olref[req][min]{lem:less-nsucc} نحصل على
(ب) $\lnot(\OLId{latin-ordinary}{x}<\num{\OLId{latin-ordinary}{n}})$؛ ومن
\olref[req][min]{lem:trichotomy} نحصل على $\num{\OLId{latin-ordinary}{n}}<\OLId{latin-ordinary}{x}$.
ونضم إلى $\Th{Q}\Proves\ORefut[\Th{T}](\num{\OLId{latin-ordinary}{n}},\gn{!R_{\Th{T}}})$
ما سبق، فيكون
$\num{\OLId{latin-ordinary}{n}}<\OLId{latin-ordinary}{x}\land\ORefut[\Th{T}](\num{\OLId{latin-ordinary}{n}},\gn{!R_{\Th{T}}})$،
ويلزم عنه
$\lexists[\OLId{latin-ordinary}{z}][(\OLId{latin-ordinary}{z}<\OLId{latin-ordinary}{x}\land\ORefut[\Th{T}](\OLId{latin-ordinary}{z},\gn{!R_{\Th{T}}}))]$.
ولاعتباطية $\OLId{latin-ordinary}{x}$ يتم المطلوب:
\[
\lforall[\OLId{latin-ordinary}{x}][(\OPrf[\Th{T}](\OLId{latin-ordinary}{x},\gn{!R_{\Th{T}}}) \lif
  \lexists[\OLId{latin-ordinary}{z}][(\OLId{latin-ordinary}{z} < \OLId{latin-ordinary}{x} \land \ORefut[\Th{T}](\OLId{latin-ordinary}{z},\gn{!R_{\Th{T}}}))])].
\]
\end{proof}

\begin{prob}
نقول إن مجموعتي الأعداد الطبيعية $\OLId{latin-looped}{A}$ و$\OLId{latin-looped}{B}$ \emph{غير قابلتين للفصل
حسابيًا} إذا امتنع وجود مجموعة $\OLId{latin-looped}{X}$ !!{decidable} تحقق
$\OLId{latin-looped}{A}\subseteq \OLId{latin-looped}{X}$ و$\OLId{latin-looped}{B}\subseteq\Complement{\OLId{latin-looped}{X}}$؛ و$\Complement{\OLId{latin-looped}{X}}$
هي المجموعة $\Nat\setminus \OLId{latin-looped}{X}$، متممة $\OLId{latin-looped}{X}$.
لتكن $\Th{T}$ امتدادًا متسقًا و!!{axiomatizable} لـ$\Th{Q}$.
واجعل $\OLId{latin-looped}{A}$ مجموعة أعداد غودل للـ!!{sentence}s التي تثبت في
$\Th{T}$، و$\OLId{latin-looped}{B}$ مجموعة أعداد غودل للجمل التي تدحض في $\Th{T}$.
أثبت امتناع الفصل الحسابي بين $\OLId{latin-looped}{A}$ و$\OLId{latin-looped}{B}$.
\end{prob}

\end{document}
